\chapter{اللقاء الثاني: قراءة في مقدمة تفسير الطبري}

\section{مدخل منهجي للبدء في تفسير الطبري}
السلام عليكم ورحمة الله وبركاته. سنبدأ الآن في دراسة كتاب «تفسير الطبري» بإذن الله تبارك وتعالى، وسيكون البدء بقراءة مقدمة التفسير. وكنت أحب أن أقرأها من الكتاب نفسه، ولكني رأيت أن الدكتور \rawi{مساعد الطيار} قد شرح هذه المقدمة، فنقرأ الشرح سوياً لنربح في طريقنا قراءة كتابين، وبعد أن ننتهي من المقدمة إن شاء الله سندخل في قراءة التفسير.

وقد رفعت لكم كتاب «شرح مقدمة تفسير الطبري» للدكتور \rawi{مساعد الطيار}، وسنبدأ القراءة من الصفحة السابعة والسبعين، وهذه المقدمة هي بالضبط مقدمة تفسير \rawi{الطبري} في كتابه، ولكن الشيخ قدم لها بقرابة سبعين صفحة تكلم فيها عن ترجمة \rawi{الطبري} رحمه الله، وعن أهمية كتابه، وطبعاته، ومقدمات عامة تفيد الطالب جداً، وسيكون من التكليفات المفيدة أن تقرأوا ترجمة \rawi{الطبري} من كتاب «سير أعلام النبلاء» أو غيره، وأن تقرأوا هذه المقدمة الماتعة.

\begin{fawaid}[أهمية سورتي الفاتحة والبقرة لاكتشاف المنهج]
تفسير \rawi{الطبري} لسورتي الفاتحة والبقرة هو أهم ما يُستكشف به منهج المؤلف، فلو أن طالباً لم تتيسر له قراءة الكتاب كاملاً واكتفى بهاتين السورتين، لاستكشف أهم ما عند \rawi{الطبري} في علم التفسير؛ لأنه توسع فيهما توسعاً كبيراً جداً، ولم يستمر على هذا التوسع لاحقاً، بل كان يعزو إليهما قائلاً: «سبق أن قلنا ذلك». ولذا أنصح من يشعر بطول الدورة أن يبقى معنا على الأقل حتى ننهي سورة البقرة.
\end{fawaid}

\section{أبواب مقدمة الطبري}
حاول الدكتور \rawi{مساعد الطيار} حفظه الله أن يذكر الأبواب التي تضمنتها مقدمة \rawi{الطبري} لتكونوا على دراية بالموضوعات التي تناولها في مقدمة تفسيره، وهي كالتالي:

\begin{enumerate}
    \item القول في البيان على اتفاق معاني آي القرآن ومعاني منطق من نزل بلسانه القرآن، ومن وجّه البيان والدلالة على أن ذلك من الله تعالى ذكره هو الحكمة البالغة. مع الإبانة عن فضل المعنى الذي به باين القرآن سائر الكلام.
    \item القول في البيان عن الأحرف التي اتفقت فيها ألفاظ العرب وألفاظ غيرها من بعض الأمم.
    \item القول في اللغة التي نزل بها القرآن من لغات العرب.
    \item القول في البيان عن معنى قول رسول الله صلى الله عليه وسلم: \begin{hadith}أنزل القرآن من سبعة أبواب الجنة\end{hadith} وذكر الأخبار الواردة بذلك.
    \item القول في الوجوه التي من قِبَلِها يُوصَل إلى معرفة تأويل القرآن.
    \item ذكر بعض الأخبار التي رويت بالنهي عن القول في تأويل القرآن بالرأي.
    \item ذكر الأخبار التي رويت في الحض على علم تفسير القرآن، ومن كان يفسره من الصحابة.
    \item ذكر الأخبار التي عُرِضَ في تأويلها منكرو القول في تأويل القرآن.
    \item ذكر فيمن كان من قدماء المفسرين محموداً علمه بالتفسير، ومن كان منهم مذموماً علمه به.
\end{enumerate}

\separator

\section{خطبة الكتاب}
\noindent هذا الموضع هو من أهم مواضع اللقاء؛ لأن الشيخ يقرأ خطبة \rawi{الطبري} ثم ينتقل منها إلى استنباط معالم منهجه في التفسير وعلوم القرآن.

\isnad{قُرِئ على} \rawi{أبي جعفر محمد بن جرير الطبري} \isnad{في سنة ست وثلاثمائة قال:}

«الحمد لله الذي حجبت الألباب بدائع حكمه، وخاصمت العقول لطائف حججه، وقطعت عذر الملحدين عجائب صنعه، وهتف في أسماع العالمين ألسن أدلته، شاهدة أنه الله الذي لا إله إلا هو، الذي لا عدل له معادل، ولا مثل له مماثل، ولا شريك له مظاهر، ولا ولد له ولا والد، ولم يكن له صاحبة ولا كفواً أحد. وإنه الجبار الذي خضعت لجبروته الجبابرة، والعزيز الذي ذلت لعزته الملوك الأعزة، وخشعت لمهابة سطوته... وأذعن له جميع الخلق بالطاعة طوعاً وكرهاً. كما قال الله عز وجل:
\begin{ayah}
وَلِلَّهِ يَسْجُدُ مَن فِي السَّمَاوَاتِ وَالْأَرْضِ طَوْعًا وَكَرْهًا وَظِلَالُهُم بِالْغُدُوِّ وَالْآصَالِ
\end{ayah}
وكل موجود إلى وحدانيته داعٍ، وكل محسوس إلى ربوبيته بما وسمهم به من آثار من نقص وزيادة وعجز وحاجة، وتصرف في عاهات عارضة ومقارنة أحداث؛ لتكون له الحجة البالغة. 
ثم أردف ما شهدت به من ذلك أدلته، وأكد ما استنارت في القلوب منه بهجته برسُل ابتعثهم إلى من يشاء من عباده دعاة إلى ما اتضحت لديهم صحته؛ لئلا يكون للناس على الله حجة بعد رسل. وليذكر أولو النهى والحلم، فأمدهم بعونه، وأبانهم من سائر خلقه بما دل به على صدقهم من الأدلة، وأيدهم به من الحجج البالغة لئلا يقول القائل فيهم:
\begin{ayah}
مَا هَٰذَا إِلَّا بَشَرٌ مِّثْلُكُمْ يَأْكُلُ مِمَّا تَأْكُلُونَ مِنْهُ وَيَشْرَبُ مِمَّا تَشْرَبُونَ * وَلَئِنْ أَطَعْتُم بَشَرًا مِّثْلَكُمْ إِنَّكُمْ إِذًا لَّخَاسِرُونَ
\end{ayah}
جعلهم سفراء بينه وبين خلقه وأمناءه على وحيه، واختصهم بفضله واصطفاهم برسالته، ثم جعلهم فيما خصهم به من مواهب، ومنّ به عليهم من كراماته مراتب مختلفة ومنازل متفرقة، ورفع بعضهم فوق بعض درجات، وتفاضلات متباينات... وفضل نبينا محمداً صلى الله عليه وسلم من الدرجات العليا، ومن المراتب العظمى، فحباه من أقسام كرامته بالقسم الأفضل، وخصاه من درجات النبوة بالحظ الأجزل، ومن الأتباع والأصحاب الأوفر، وابتعثه بالدعوة التامة والرسالة العامة... حتى أظهر به الدين وأوضح به السبيل، وأنهج به معالم الحق ومحق به منار الشرك، وزهق به الباطل واضمحل به الضلال وخُدِع الشيطان وعبادة الأصنام والأوثان.

أما بعد، فإن مما خص الله به أمة نبينا محمد صلى الله عليه وسلم من الفضيلة وشرفهم به على سائر الأمم من المنازل الرفيعة، وحباهم به من الكرامة السنية: حفظه ما حفظ جل ذكره وتقدست أسماؤه عليهم من وحيه وتنزيله الذي جعله على حقيقة نبوة نبيهم صلى الله عليه وسلم، وأعلى ما خصه به من الكرامة علامة واضحة وحجة بالغة، أبانه به من كل كاذب ومفتر، وفصل به بينهم وبين كل جاحد وملحد... الذي لو اجتمع جميع من بين أقطارها من جنها وإنسها وصغيرها وكبيرها على أن يأتوا بسورة لمثله لم يأتوا بمثله ولو كان بعضهم لبعض ظهيراً. قد جعله لهم في دجى الظلم نوراً ساطعاً، وفي سدف الشبه شهاباً لامعاً...
\begin{ayah}
يَهْدِي بِهِ اللَّهُ مَنِ اتَّبَعَ رِضْوَانَهُ سُبُلَ السَّلَامِ وَيُخْرِجُهُم مِّنَ الظُّلُمَاتِ إِلَى النُّورِ بِإِذْنِهِ وَيَهْدِيهِمْ إِلَىٰ صِرَاطٍ مُّسْتَقِيمٍ
\end{ayah}
... فهو موئلهم الذي إليه عند الاختلاف يئلون، ومعقلهم الذي إليه في النوازل يعقلون، وحصنهم الذي به من وساوس الشيطان يتحصنون، وحكمة ربهم التي إليها يحتكمون... 

اللهم فوفقنا لإصابة صواب القول في محكمه ومتشابهه وحلاله وحرامه وعامه وخاصه ومجمله ومفسره وناسخه ومنسوخه وظاهره وباطنه وتأويل آيه وتفسير مشكله، والتمسك به والاعتصام بمحكمه، والثبات على التسليم لمتشابهه... واعلموا عباد الله رحمكم الله أن أحق ما صرفت إلى علمه العناية، وبلغت في معرفته الغاية، ما كان لله في العلم به رضا... وتنزيله الذي لا مرية فيه الفائز بجزيل الذخر وسنة الأجر تاليه:
\begin{ayah}
لَّا يَأْتِيهِ الْبَاطِلُ مِن بَيْنِ يَدَيْهِ وَلَا مِنْ خَلْفِهِ ۖ تَنزِيلٌ مِّنْ حَكِيمٍ حَمِيدٍ
\end{ayah}

ونحن في شرح تأويله وبيان ما فيه من معانيه منشئون إن شاء الله ذلك كتاباً مستوعباً لكل ما بالناس إليه الحاجة من علمه، جامعاً ومن سائر الكتب غيره في ذلك كافياً. ومخبرون في كل ذلك بمنتهى إلينا من اتفاق فيما اتفقت عليه الأمة واختلافها فيما اختلفت فيه منه. ومبينو علل كل مذهب من مذاهبهم... باوجز ما أمكن من الإيجاز في ذلك وأخصر ما أمكن من الاختصار فيه».

\begin{fawaid}[أهمية قراءة النصوص مضبوطة بالشكل]
الأفضل عند قراءة كلام \rawi{الطبري} أن نرجع إلى النسخة الأصلية من الكتاب لأنها مضبوطة بالتشكيل؛ فكلام \rawi{الطبري} دقيق ويحتاج إلى تشكيل كثير ليُفهم على وجهه الصحيح كمعرفة الفاعل والمبني للمجهول في قوله: (وخُدِعَ الشيطانُ).
\end{fawaid}

\separator

\section{الشرح والتعليق على خطبة الكتاب}
استهل \rawi{الطبري} رحمه الله مقدمة تفسيره بخطبة أثنى فيها على الله تعالى، ومهد بها للدخول في موضوع كتابه.

\begin{fawaid}[منهجية شرح كتب التراث]
نلاحظ أمراً مهماً؛ وهو أن الشيخ الشارح لم يشرح كل لفظة في المقدمة، وهذا هو الصواب. فكثير من الناس إذا جاء يشرح كلام أهل العلم يشرح كل لفظ فيه حتى البَيِّن. الأمور الميسرة التي تُفهم بذاتها لا تحتاج لتعليق، وإنما العناية والتفصيل تكون للأمور التي تحتاج إلى الفحص والدراسة.
\end{fawaid}

توقف الشارح عند قول \rawi{الطبري}: «اللهم فوفقنا لإصابة صواب القول في محكمه ومتشابهه وحلاله وحرامه وعامه وخاصه ومجمله وناسخه ومنسوخه وظاهره وباطنه وتأويل آيه وتفسير مشكله».
وهنا ندرك أن علم تفسير القرآن يندرج تحت دائرة أوسع وهي "علوم القرآن"؛ وهي العلوم المتعلقة بكتاب الله كنزول القرآن، وكتابة المصحف، وعلم القراءات، والمكي والمدني وغيرها.

وقد طرح الشارح عدة تساؤلات حول هذه الفقرة:

\subsection{أولاً: هل هذه العلوم متقابلة؟}
نلاحظ أن (المحكم) مقابل للمتشابه، و(الحلال) مقابل للحرام، و(العام) مقابل للخاص، وهكذا. ولكن الموطن الأخير «وتأويل آية، وتفسير مشكله» فيه إشكال؛ إذ يقتضي مقابلة (التأويل) بـ (التفسير)، مع أن مصطلح (التفسير) نادر الدوران في كتاب \rawi{الطبري} مقارنة بمصطلح (التأويل). 

والراجح أن هذه العبارة من باب تنويع العبارة وعطف الصفات لا المتقابلات، لأنه لا يظهر في كتابه أثر للتفريق بين التفسير والتأويل.

\subsection{ثانياً: هل ذكره لهذه العلوم على سبيل الحصر أم التمثيل؟}
لا يمكن أن نقول إنه أراد هنا الاستيعاب والحصر، إنما هو مجرد تمثيل. ومما يجب التنبه له أن علوم القرآن نوعان:
\begin{itemize}
    \item علوم تضمنها القرآن (موجودة في النص نفسه): مثل العام والخاص.
    \item علوم محتفة بالنص (تأخذ صبغة تاريخية): مثل المكي والمدني وأسباب النزول.
\end{itemize}
والذي يبدو أن \rawi{الطبري} يتكلم هنا عن النوع الأول؛ أي صفات النص نفسه والمحتملات الواردة على الآية، كأن يكون النص عاماً أو خاصاً، محكماً أو منسوخاً.

\subsection{ثالثاً: ما مصادر الطبري في ذكر هذه العلوم؟}
يمكن الإشارة إلى أربعة مصادر يُظن أن \rawi{الطبري} استقى منها هذه التقسيمات، والمراد هنا تقريب مأخذها لا الجزم بأن هذه القائمة على سبيل الحصر:
\begin{enumerate}
    \item \textbf{القرآن الكريم:} إذ فيه ذكر بعض هذه الأنواع، كقوله تعالى: \begin{ayah}مِنْهُ آيَاتٌ مُّحْكَمَاتٌ هُنَّ أُمُّ الْكِتَابِ وَأُخَرُ مُتَشَابِهَاتٌ\end{ayah}.
    \item \textbf{قول النبي صلى الله عليه وسلم:} كما روى \rawi{الطبري}: \begin{hadith}كان الكتاب الأول نزل من باب واحد وعلى حرف واحد، ونزل القرآن من سبعة أبواب وعلى سبعة أحرف: زجر، وأمر، وحلال، وحرام، ومحكم، ومتشابه، وأمثال. فأحلوا حلاله، وحرموا حرامه، وافعلوا ما أمرتم به، وانتهوا عما نهيتم عنه، واعتبروا بأمثاله، واعملوا بمحكمه، وآمنوا بمتشابهه\end{hadith}.
    \item \textbf{ما ورد عن السلف:} كأثر \rawi{ابن مسعود}: «والله الذي لا إله غيره ما أنزلت سورة من كتاب الله إلا وأنا أعلم أين نزلت...»، وكاستدلاله على العام والخاص بالجمع بين الآيتين: \begin{ayah}وَأُولَاتُ الْأَحْمَالِ أَجَلُهُنَّ أَن يَضَعْنَ حَمْلَهُنَّ\end{ayah} و \begin{ayah}وَالْمُطَلَّقَاتُ يَتَرَبَّصْنَ بِأَنفُسِهِنَّ ثَلَاثَةَ قُرُوءٍ\end{ayah}.
    \item \textbf{كتاب «الرسالة» للإمام \rawi{الشافعي}:} فقد تكلم \rawi{الشافعي} رحمه الله عن كثير من هذه العلوم كالعام والخاص والناسخ والمنسوخ، والطبري قد استفاد منه كثيراً وظهر أثره جلياً في مقدمته.
\end{enumerate}

\begin{fawaid}[بناء الخلاصات العلمية]
الخلاصة التي يذكرها العالم المحقق هي نتاج علمه بمن سبقه وجمعه وتحريره. فمعرفة مصدر قول العالم تعين جداً على فهم القول، ومعرفة كيف بنى العالم اختياراته.
\end{fawaid}

\separator

\section{معالم منهج الطبري في التفسير}
في قوله: «ونحن في شرح تأويله وبيان ما فيه من معانيه منشئون إن شاء الله ذلك كتاباً مستوعباً...»، تتجلى أهم فقرة في المقدمة لبيان منهج \rawi{الطبري}، ويمكن تلخيص معالمها في النقاط التالية:

\subsection{أولاً: الاستيعاب وتبيين معاني القرآن}
إطار عمل \rawi{الطبري} هو (تبيين معاني القرآن الكريم وتأويله). وهو ما صرح به، فجاء تفسيره خِلْواً من كثير من الحشو والاستطراد الذي لا يخدم هذا الغرض (كالتوسع في الإعراب والقراءات والأحكام الفقهية إلا ما مست إليه الحاجة لفهم المعنى). 

قال واصفاً كتابه: «كتاباً مستوعباً لكل ما بالناس إليه الحاجة من علمه جامعاً، ومن سائر الكتب غيره في ذلك كافياً». وهذا الاستيعاب يجعله كتاباً تأسيسياً لا يمكن لطالب التفسير أن يستغني عنه أو يعوضه بغيره. 

\begin{fawaid}[أهمية معرفة خاصة الكتاب وغايته]
لا يصح محاكمة كتاب لهدف لم يقصده وغاية لم يتطلبها، فلا يقال: "تفسير فلان أشمل لبلاغة القرآن أو لوجوه إعرابه من تفسير الطبري". لأن \rawi{الطبري} التزم غرضاً واحداً وهو "بيان المعاني"، ولم يجعله موسوعة نحوية أو فقهية، وهذه هي خاصة الكتاب التي يجب أن تُفهم لكي يُنتفع منه على الوجه الصحيح.
\end{fawaid}

\subsection{ثانياً: إيراد الاتفاق والاختلاف}
قال \rawi{الطبري}: «ومخبرون في كل ذلك بما انتهى إلينا من اتفاق الحجة فيما اتفقت عليه الأمة، واختلافها فيما اختلفت فيه منه». 

كتابه مبني على تصوير المسألة بذكر موضع الإجماع ثم إيراد الاختلاف، والمقصود باستيعابه هو استيعاب (الأقوال والمذاهب) وليس بالضرورة استيعاب كل (المرويات والأسانيد) المكررة في ذات المعنى، ومصطلح "الحجة" عنده ينصرف إلى الإجماع عند طبقات الصحابة والتابعين وأتباعهم.

\subsection{ثالثاً: توجيه الأقوال وتبيين العلل}
قال: «ومبينو علل كل مذهب من مذاهبهم». وهذه ميزة كبرى؛ فهو لا يكتفي بذكر القول، بل يذكر العلة والدليل والأصل الذي بُني عليه هذا القول، سواء قَبِلَه أو رده. وتوجيه الأقوال علم عزيز يوسّع المدارك في فهم آراء السلف.

\subsection{رابعاً: الترجيح بين الأقوال المختلفة}
قال: «وموضحوا الصحيح لدينا من ذلك». فهو يرجح بعد الجمع والتصنيف والنقد، ويبني ترجيحه على أدلة وقرائن وموجبات واضحة، ولا يلقي الكلام على عواهنه. ومن العلماء الذين استفادوا جداً من ترجيحات \rawi{الطبري} وقواعده الإمام \rawi{أبو جعفر النحاس}.

\subsection{خامساً: الإيجاز والاختصار}
قال: «بأوجز ما أمكن من الإيجاز في ذلك، وأخصر ما أمكن من الاختصار فيه». رغم غزارة علم الكتاب وحجمه، إلا أن \rawi{الطبري} اختصره قدر الإمكان، مقدماً أصح معلومة بأوجز عبارة، وهذا يدل على قدرته الفائقة على ضبط منهجه، كما فعل في كتابه «تاريخ الأمم والملوك».

\separator

\section{موازنة تفسير الطبري بالتفاسير السابقة له}
لإدراك القفزة النوعية التي أحدثها \rawi{الطبري}، نجد أن الكتب التي جمعت قبله (مثل «تفسير مقاتل بن سليمان» و«تفسير يحيى بن سلام») كانت تعتمد إما على الاختيار أو الجمع المجرد دون نقد أو تعليل أو ترجيح مستند إلى قواعد. 

وقد ظن \rawi{الفاضل بن عاشور} في كتابه «التفسير ورجاله» أن \rawi{الطبري} مدين لمنهجه لـ \rawi{يحيى بن سلام}، وأنه يمثل "الحلقة المفقودة"، ولكن هذا مردود؛ فـ \rawi{الطبري} لم يطلع على تفسيره يقيناً لأنه صُنف في القيروان ولم يظهر بالمشرق في وقته، بل إن \rawi{الطبري} أبدع هذه الطريقة النقدية والتحليلية بنفسه، فكان العلم بعده ليس كالعلم قبله.

\section{أهمية علوم العربية لتفسير القرآن}
ختم \rawi{الطبري} مقدمته بالتأكيد على أهمية لغة العرب قائلاً: «وذلك البيان عما في آي القرآن من المعاني التي من قِبَلِها يدخل اللبس على من لم يُعانِ رياضة العلوم العربية ولم تستحكم معرفته بتصاريف وجوه منطق الألسن السليقية الطبيعية».

وهذه الجملة تشبه صنيع الإمام \rawi{الشافعي} في «الرسالة»، مما يؤكد قوة الصلة بين طريقة \rawi{الطبري} وطريقة \rawi{الشافعي} في تأسيس النظر. فبما أن القرآن نزل بلسان العرب، فإن مدار دراسة فهم الآيات والتفسير يرتكز على العناية بعلوم العربية، ومن جهلها دخلت عليه الشبهات واللبس.

\separator

بهذا نكون قد مررنا على أهم معالم خطبة \rawi{الطبري} للكتاب، وسنكمل غداً إن شاء الله استعراض باقي المقدمة، قبل الشروع في التفسير العملي لسورة الفاتحة.

والسلام عليكم ورحمة الله وبركاته.
